\section{Реализация, тестирование и демонстрация гибридной интеллектуальной системы}

Разработка программного комплекса автономной видеорегистрации и анализа шахматных партий выполнена в соответствии с архитектурными решениями, математическими моделями и онтологической структурой предметной области, описанными в разделе~2. Созданная система объединяет низкоуровневые алгоритмы компьютерного зрения, нейросетевые модели классификации клеток и целостного анализа доски, асинхронный конечно-автоматный фильтр помех, логико-символьный верификатор правил ФИДЕ, транзакционный репозиторий с криптографическим хеш-цепочечным связыванием, веб-панель оператора на базе \texttt{FastAPI} и службу шахматного анализа на базе UCI-движка \texttt{Stockfish~16}.

\subsection{Выбор средств и технологий реализации}

Выбор технологического стека обусловлен требованиями надёжности работы на стандартном центральном процессоре, строгим контролем задержек цикла захвата ($\le 35$~мс на кадр в фазе детекции движения и стабилизации), строгой типизацией межмодульных контрактов и переносимостью между ОС Linux и Docker-контейнерами:

\begin{enumerate}
    \item \textbf{Язык программирования Python 3.10--3.14}. Выбран в качестве базового языка интеграции благодаря развитой экосистеме библиотек машинного обучения и компьютерного зрения, поддержке статической типизации (\texttt{typing}), строгой валидации контрактов данных через \texttt{pydantic~v2} и примитивам многопоточности (\texttt{threading}, \texttt{queue}).
    \item \textbf{Фреймворки глубокого обучения и инференса (PyTorch~2.2+~\cite{Paszke2019} и ONNX Runtime~1.17+~\cite{ONNXRuntime2021})}. Обучение и квантование моделей выполнено на базе PyTorch, а промышленный CPU-инференс переведён в формат ONNX. В гибридном конвейере используются как компактный свёрточный поклеточный классификатор (\texttt{EfficientNet-B0}~\cite{Wolflein2021}), так и целостный визуальный трансформер всей доски \texttt{ChessQ-Lite (ViT-S14 INT8)}~\cite{Dosovitskiy2021} совместно с патч-эмбеддингами \texttt{DINOv2-Small INT8}~\cite{Oquab2024} для инвариантного к освещению сравнения состояний клеток.
    \item \textbf{Библиотека компьютерного зрения OpenCV 4.9+~\cite{Bradski2000}}. Используется для захвата и нормализации кадров, цветовых преобразований, гауссова сглаживания, детектирования границ оператором Кэнни, вероятностного преобразования Хафа, 1D-проективной и RANSAC-оценки~\cite{Fischler1981} четырёх углов доски, расчёта матрицы перспективной гомографии $H$~\cite{Szeliski2022} и геометрического выпрямления доски к размеру $640 \times 640$ пикселей.
    \item \textbf{Комплекс отслеживания рук (MediaPipe Hand Landmarker~\cite{Lugaresi2019})}. Обеспечивает пространственную 21-точечную детекцию скелета кисти руки игрока над игровой зоной. Для сред без модели или при отсутствии явной кисти дополнительно действует цветово-силуэтный детектор \texttt{ColorSilhouetteHandDetector} в пространстве YCrCb.
    \item \textbf{Логико-символьное ядро шахматных правил (python-chess 1.10+~\cite{Fiekas2024})}. Предоставляет строгое верифицированное представление правил Международной шахматной федерации~\cite{FIDE2023}: генерацию легальных ходов, распознавание шахов, матов, патов, троекратного повторения, правила 50 ходов и сериализацию в форматы FEN и PGN.
    \item \textbf{Подсистема компьютерного анализа (Stockfish 16 UCI~\cite{Romstad2024})}. Открытый шахматный движок с нейросетевой оценочной функцией NNUE. Интегрирован через асинхронный протокол UCI с очередью задач и эпохальной отменой устаревших запросов при сбросе или откате партии.
    \item \textbf{Веб-стекинг и контейнеризация (FastAPI, Uvicorn, WebSocket, MJPEG, Docker Compose)}. Обеспечивает кроссплатформенную операторскую панель в браузере, трансляцию видеопотока \texttt{MJPEG}, реактивное обновление состояния партии через \texttt{WebSocket}, запись видеоархива и офлайн-транскрибацию видеофайлов.
    \item \textbf{Средства тестирования (pytest 9.2+)}. Автоматизированный запуск модульных, интеграционных, отказоустойчивых (fault-injection), веб-интерфейсных и сквозных тестов.
\end{enumerate}

Сводное сопоставление выбранных технологий и реализуемых ими архитектурных компонентов приведено в таблице~\ref{tab:tech-stack}.

\begin{table}[htbp]
\centering
\caption{Фактический технологический стек программной реализации}
\label{tab:tech-stack}
\begin{tabular}{|>{\raggedright\arraybackslash}p{0.27\textwidth}|>{\raggedright\arraybackslash}p{0.21\textwidth}|>{\raggedright\arraybackslash}p{0.44\textwidth}|}
\hline
\textbf{Технология / Инструмент} & \textbf{Версия} & \textbf{Назначение в архитектуре} \\
\hline
Python & 3.10 -- 3.14 & Единая среда интеграции компонентов, бизнес-логика, CLI. \\
\hline
Pydantic & 2.6+ & Строго типизированные контракты данных (\texttt{Frame}, \texttt{MoveEvent}, \texttt{BoardGeometry}). \\
\hline
OpenCV & 4.9+ & Проективная геометрия, фильтрация движения, Hough/Huber/RANSAC, HUD. \\
\hline
PyTorch / ONNX Runtime & 2.2+ / 1.17+ & Обучение и CPU-инференс \texttt{EfficientNet-B0}, \texttt{ChessQ-Lite ViT} и \texttt{DINOv2}. \\
\hline
MediaPipe & 0.10+ & Пространственная детекция кисти руки и расчёт окклюзии 64 полей. \\
\hline
python-chess & 1.10+ & Авторитетный символьный решатель, правила ФИДЕ, FEN/PGN. \\
\hline
Stockfish & 16 (UCI) & Фоновый шахматный анализ позиций, шкала перевеса и классификация ошибок. \\
\hline
FastAPI / WebSocket & 0.110+ & Веб-сервер панели оператора, \texttt{MJPEG}-поток, REST API и синхронизация состояния. \\
\hline
Docker Compose & v2 & Изолированное развёртывание контейнеров \texttt{chess-recorder-web} и \texttt{gui}. \\
\hline
Pytest & 9.2+ & Модульное, интеграционное, веб-API и сквозное тестирование (285 тестов). \\
\hline
\end{tabular}
\end{table}

\subsection{Структура программной реализации}

Исходный код системы структурирован по модульному принципу в пакете \texttt{src/} с чётким разделением контуров:

\begin{itemize}
    \item \texttt{src/core/} --- контракты данных (\texttt{datatypes.py}), конфигурация (\texttt{config.py}), логирование и исключения;
    \item \texttt{src/vision/} --- захват видео (\texttt{VideoAdapter}), локализация и трекинг углов доски (\texttt{BoardLocalizer}, \texttt{CornerRefiner}, \texttt{BoardGeometryTracker}), гомография, трекинг занятости зоны основания фигур (\texttt{Base-Core}), оптическое движение и детекция рук;
    \item \texttt{src/ml/} --- поклеточный классификатор (\texttt{ONNXPieceClassifier}), целостный трансформер всей доски (\texttt{WholeBoardChessQClassifier}), детектор изменений по патч-эмбеддингам \texttt{DINOv2PatchChangeDetector}, авторазметчик и пайплайн обучения;
    \item \texttt{src/reasoning/} --- конечный автомат \texttt{GameFSM}, генератор легальных кандидатов, многофакторный скорер, шлюз принятия решений \texttt{DecisionGate} и авторитетный менеджер \texttt{GameStateManager};
    \item \texttt{src/persistence/} --- транзакционный журнал \texttt{events.jsonl} с SHA-256 хеш-цепочками, сохранение фотодоказательств и модуль восстановления \texttt{RecoveryManager};
    \item \texttt{src/analysis/} --- асинхронный супервизор процесса Stockfish, потокобезопасная очередь задач с эпохальной отменой и оценщик ходов по сантипешкам;
    \item \texttt{src/ui/} --- веб-сервер \texttt{src/ui/web/} (\texttt{FastAPI}, \texttt{MJPEGVideoStreamer}, \texttt{WebSocketManager}, видеосервисы записи и офлайн-анализа, SPA-интерфейс), оверлей \texttt{VisualHUDComposer}, десктопный \texttt{InteractiveDashboard} и точка входа CLI (\texttt{src/ui/cli.py}).
\end{itemize}

Детальная спецификация программных модулей, реализованных классов и их функций представлена в таблице~\ref{tab:code-structure}.

\begin{longtable}{|>{\raggedright\arraybackslash}p{0.23\textwidth}|>{\raggedright\arraybackslash}p{0.27\textwidth}|>{\raggedright\arraybackslash}p{0.42\textwidth}|}
\caption{Структурная декомпозиция программной реализации}\label{tab:code-structure} \\
\hline
\textbf{Модуль} & \textbf{Ключевые классы} & \textbf{Назначение и функционал} \\
\hline
\endfirsthead
\hline
\textbf{Модуль} & \textbf{Ключевые классы} & \textbf{Назначение и функционал} \\
\hline
\endhead
\texttt{src.core.datatypes} & \texttt{Frame}, \texttt{BoardGeometry}, \texttt{CellObservation}, \texttt{MoveEvent}, \texttt{PipelineState} & Строго типизированные контракты Pydantic для передачи наблюдений и событий по конвейеру. \\
\hline
\texttt{src.core.config} & \texttt{AppConfig}, \texttt{VisionConfig}, \texttt{TemporalConfig}, \texttt{DecisionConfig} & Иерархическая загрузка параметров из YAML с переопределением через переменные окружения. \\
\hline
\texttt{src.vision.board\_localizer} & \texttt{CoarseToFineLocalizer}, \texttt{BoardGeometryTracker} & Двухстадийная детекция углов (YOLOv8 + Canny/Hough/Huber/RANSAC + 1D-проективная проверка) и EMA-трекинг дрейфа. \\
\hline
\texttt{src.vision.homography} & \texttt{GeometryTransformer} & Расчёт матрицы $H \in \mathbb{R}^{3 \times 3}$, нормализация кадра к $640 \times 640$ пикселей, проверка выпуклости и обратных ошибок. \\
\hline
\texttt{src.vision.occupancy\_tracker} & \texttt{BoardOccupancyTracker} & Инвариантная к яркости и теням оценка занятости клеток в зоне основания фигуры (\textit{Base-Core}). \\
\hline
\texttt{src.vision.motion\_detector}, \texttt{hand\_detector} & \texttt{MotionDetector}, \texttt{HandDetector}, \texttt{ColorSilhouetteHandDetector} & Подавление глобальных скачков яркости, межкадровый анализ и 21-точечная маска перекрытия рукой. \\
\hline
\texttt{src.ml.piece\_classifier} & \texttt{ONNXPieceClassifier}, \texttt{WholeBoardChessQClassifier}, \texttt{DINOv2PatchChangeDetector} & Пакетный инференс 64 клеток, целостное распознавание всей доски \texttt{ChessQ-Lite ViT-S14} и сравнение патчей \texttt{DINOv2}. \\
\hline
\texttt{src.reasoning.fsm} & \texttt{GameFSM} & Конечный автомат фильтрации помех (\texttt{PAUSED}, \texttt{IDLE}, \texttt{CALIBRATING}, \texttt{STABLE}, \texttt{MOVING}, \texttt{SETTLING}, \texttt{VERIFYING}, \texttt{UNCERTAIN}). \\
\hline
\texttt{src.reasoning.*} & \texttt{LegalMoveCandidateGenerator}, \texttt{MoveScorer}, \texttt{DecisionGate}, \texttt{GameStateManager} & Топологическая фильтрация легальных ходов ФИДЕ, скоринг $S(m)$, проверка запаса $\Delta$ и авторитетное обновление позиции. \\
\hline
\texttt{src.persistence.*} & \texttt{JSONLEventRepository}, \texttt{EvidenceWriter}, \texttt{RecoveryManager} & Атомарный аудит \texttt{events.jsonl} с SHA-256 цепочкой (\texttt{move}, \texttt{rollback}, \texttt{illegal\_attempt}) и усечением битого хвоста. \\
\hline
\texttt{src.analysis.*} & \texttt{StockfishWorker}, \texttt{AsyncEngineService}, \texttt{MoveEvaluator} & Неблокирующий UCI-анализ с эпохальной отменой устаревших задач и оценкой потерь в сантипешках. \\
\hline
\texttt{src.orchestrator} & \texttt{SystemOrchestrator} & Сквозной дирижёр конвейера: управление сессиями, калибровкой, гибридной верификацией и публикацией состояний. \\
\hline
\texttt{src.ui.web.*} & \texttt{app}, \texttt{MJPEGVideoStreamer}, \texttt{WebSocketManager}, \texttt{VideoTranscriptionManager} & FastAPI-бэкенд, потоковая передача MJPEG, двусторонняя синхронизация WS, выбор V4L2/видео и офлайн-анализ файлов. \\
\hline
\texttt{src.ui.visualizer}, \texttt{dashboard}, \texttt{cli} & \texttt{VisualHUDComposer}, \texttt{InteractiveDashboard}, \texttt{main} & Графический оверлей, локальный OpenCV HUD и CLI с подкомандами \texttt{record}, \texttt{web}, \texttt{audit}, \texttt{replay}, \texttt{benchmark}. \\
\hline
\end{longtable}

\subsection{Реализация пользовательского интерфейса}

В соответствии с требованиями академической и эксплуатационной строгости программный комплекс реализует трёхуровневый пользовательский интерфейс:
\begin{enumerate}
    \item \textbf{Интерактивная веб-панель оператора (Web Dashboard SPA)}: основное рабочее место пользователя на базе \texttt{FastAPI}, протокола \texttt{WebSocket} (\texttt{/ws/live}) и потокового видео \texttt{MJPEG} (\texttt{/stream/video}), не требующее локального X11/Wayland-сервера при запуске в Docker-контейнере.
    \item \textbf{Локальный графический экран (Visual HUD Dashboard)}: автономное окно на OpenCV HighGUI с автоматическим переходом в режим \texttt{headless} при отсутствии дисплея.
    \item \textbf{Консольный интерфейс (CLI)}: подкоманды \texttt{record}, \texttt{web}, \texttt{audit}, \texttt{replay} и \texttt{benchmark} для автоматизированного и пакетного применения.
\end{enumerate}

Веб-панель оператора логически разделена на верхнюю командную панель управления и двухоконную рабочую область (рисунки~\ref{fig:design-architecture} и~\ref{fig:ui-web-main}):
\begin{itemize}
    \item \textbf{Командная панель}: селектор источника видео (камеры \texttt{/dev/video*}, RTSP-потоки и загруженные видеозаписи), строго формализованные кнопки управления процессом (\textit{«Запуск / Пауза»}, \textit{«Калибровка»}, \textit{«Перекалибровка»}, \textit{«Отмена»}, \textit{«Сброс»}), менеджер сессий, инструменты записи видеопотока, запуска офлайн-анализа видеофайла и экспорта протокола в стандарт \texttt{PGN}.
    \item \textbf{Левая область (Видеоканал и геометрия)}: потоковое изображение с наложенной четырёхугольной сеткой доски и буквенно-цифровой координатной привязкой, интерактивные угловые маркеры для ручной подстройки с пересчётом координат под режим \texttt{object-fit: contain}, бейдж состояния FSM и строка телеметрии (FPS, номер кадра, номер полухода).
    \item \textbf{Правая область (Позиция, история и анализ)}: интерактивная виртуальная 2D-доска, вертикальная шкала преимущества \texttt{Stockfish} (Eval Gauge), прокручиваемая таблица ходов с возможностью пошагового ретроспективного просмотра позиции и цветовой маркировкой качества хода (\texttt{BEST}, \texttt{INACCURACY}, \texttt{MISTAKE}, \texttt{BLUNDER}).
\end{itemize}

При возникновении спорной ситуации или физической перестановки фигур, не соответствующей ни одному легальному ходу ФИДЕ, система переводит автомат в состояние \texttt{UNCERTAIN}, фиксирует инцидент \texttt{illegal\_attempt} в журнале и отображает неблокирующую сворачиваемую панель диагностики с перечнем изменённых полей, топ-кандидатами (при наличии) и полем ручного ввода хода в нотации UCI. В отличие от ранних прототипов, панель не перекрывает видеоканал и не препятствует возврату фигур на доску для автоматического продолжения партии.

На рисунках~\ref{fig:ui-web-main}--\ref{fig:ui-history-eval} предусмотрены области для размещения скриншотов экранных форм веб-панели оператора.

\begin{figure}[htbp]
\centering
\fbox{\parbox{0.92\textwidth}{\centering\vspace{3.4cm}\textit{[Место для скриншота интерфейса: главная рабочая область веб-панели оператора во время видеорегистрации партии]}\vspace{3.4cm}}}
\caption{Главная страница веб-панели оператора в процессе видеорегистрации партии}
\label{fig:ui-web-main}
\end{figure}

\begin{figure}[htbp]
\centering
\fbox{\parbox{0.92\textwidth}{\centering\vspace{3.4cm}\textit{[Место для скриншота интерфейса: режим калибровки доски с интерактивными угловыми маркерами и привязкой осей]}\vspace{3.4cm}}}
\caption{Интерактивная калибровка и перекалибровка углов шахматной доски}
\label{fig:ui-calibration}
\end{figure}

\begin{figure}[htbp]
\centering
\fbox{\parbox{0.92\textwidth}{\centering\vspace{3.4cm}\textit{[Место для скриншота интерфейса: сворачиваемая панель диагностики UNCERTAIN и фиксации нелегального хода]}\vspace{3.4cm}}}
\caption{Сворачиваемая панель разрешения неоднозначности и диагностики нелегального хода}
\label{fig:ui-uncertain}
\end{figure}

\begin{figure}[htbp]
\centering
\fbox{\parbox{0.92\textwidth}{\centering\vspace{3.4cm}\textit{[Место для скриншота интерфейса: таблица ходов, пошаговый просмотр позиции и шкала оценки Stockfish]}\vspace{3.4cm}}}
\caption{Аналитическая панель с историей ходов, шкалой Stockfish и экспортом PGN}
\label{fig:ui-history-eval}
\end{figure}

\subsection{Тестирование интеллектуальной системы}

Надёжность, детерминизм и отказоустойчивость программного комплекса подтверждены многоуровневым комплексом автоматизированного тестирования. Набор включает \textbf{285 тестов} в 38 тестовых модулях, выполняемых за $45{,}46$~с со 100\% успешностью.

Помимо функционального тестирования был проведён комплексный шестиагентный аудит кодовой базы (веб-интерфейс и API, сервисы/CLI/Docker, оркестратор и FSM, персистентность и Stockfish, компьютерное зрение и ML), по результатам которого устранены критические и архитектурные дефекты:
\begin{itemize}
    \item устранена гонка потоков при переключении видеоисточника V4L2 и нормализован пересчёт координат мыши при \texttt{object-fit: contain};
    \item введена устойчивая M-оценка Хубера (\texttt{cv2.DIST\_HUBER}), детерминированная инициализация ГПСЧ при выделении отрезков Хафа и строгий штраф внутренних подсеток доски;
    \item обеспечено корректное определение нелегальных телепортов фигур с регистрацией записи \texttt{illegal\_attempt} и их фильтрацией при формировании PGN;
    \item реализовано атомарное восстановление журнала событий при обрыве записи (усечение \texttt{r+b} через \texttt{truncate} и \texttt{fsync}), эпохальная отмена устаревших задач Stockfish и безопасный откат при закрытом движке;
    \item исправлена оценка доставленного мата (\texttt{mate\_in = 0}) без ложной классификации как грубой ошибки.
\end{itemize}

Сводные результаты автоматизированного тестирования по функциональным группам модулей представлены в таблице~\ref{tab:test-results}.

\begin{table}[htbp]
\centering
\caption{Сводные результаты тестирования подсистем комплекса}
\label{tab:test-results}
\begin{tabular}{|>{\raggedright\arraybackslash}p{0.46\textwidth}|>{\centering\arraybackslash}p{0.14\textwidth}|>{\centering\arraybackslash}p{0.16\textwidth}|>{\raggedright\arraybackslash}p{0.16\textwidth}|}
\hline
\textbf{Группа тестовых модулей} & \textbf{Кол-во тестов} & \textbf{Статус} & \textbf{Уровень проверки} \\
\hline
Контракты данных и конфигурация (\texttt{test\_datatypes}, \texttt{test\_config}) & 39 & Пройдено & Модульный \\
\hline
Компьютерное зрение и геометрия (\texttt{test\_board\_localizer}, \texttt{test\_homography}, \texttt{test\_corner\_refinement}, \texttt{test\_orientation}, \texttt{test\_motion\_detector}, \texttt{test\_hand\_detector}, \texttt{test\_roi\_extractor}, \texttt{test\_occupancy\_tracker}, \texttt{test\_video\_adapter}, \texttt{test\_diagnostics}) & 75 & Пройдено & Модульный и регрессионный \\
\hline
Нейросетевой уровень и обучение (\texttt{test\_piece\_classifier}, \texttt{test\_auto\_labeller}, \texttt{test\_dataset\_split}, \texttt{test\_training\_pipeline}, \texttt{test\_transforms}) & 28 & Пройдено & Модульный \\
\hline
Логико-символьный решатель и FSM (\texttt{test\_fsm}, \texttt{test\_candidate\_generator}, \texttt{test\_move\_scorer}, \texttt{test\_decision\_gate}, \texttt{test\_game\_state\_manager}, \texttt{test\_reasoning\_e2e}) & 39 & Пройдено & Модульный и интеграционный \\
\hline
Персистентность и анализ Stockfish (\texttt{test\_event\_repository}, \texttt{test\_evidence\_writer}, \texttt{test\_recovery\_manager}, \texttt{test\_stockfish\_worker}, \texttt{test\_engine\_service\_queue}, \texttt{test\_move\_evaluator}, \texttt{test\_persistence\_analysis\_e2e}) & 22 & Пройдено & Отказоустойчивость и UCI \\
\hline
Оркестратор, Web API, HUD и CLI (\texttt{test\_orchestrator}, \texttt{test\_orchestrator\_parallax}, \texttt{test\_whole\_board\_and\_audit}, \texttt{test\_visualizer}, \texttt{test\_cli}, \texttt{test\_web\_ui}) & 74 & Пройдено & Интеграционный и E2E \\
\hline
Эталонные системные бенчмарки (\texttt{test\_benchmark}, \texttt{test\_system\_benchmark}) & 8 & Пройдено & Сквозной системный \\
\hline
\textbf{Итого по 38 тестовым файлам} & \textbf{285} & \textbf{100\%} & \textbf{45,46 с} \\
\hline
\end{tabular}
\end{table}

\subsection{Оценка качества, эффективности и объяснимости системы}

Экспериментальная проверка производительности и качества выполнена в двух режимах: на синтетических сквозных видеопоследовательностях (\texttt{test\_system\_benchmark.py}) и непосредственно внутри рабочего контейнера \texttt{chess-recorder-web} с загруженными квантованными моделями \texttt{ONNX INT8} и движком \texttt{Stockfish~16} на центральном процессоре.

Результаты замеров ключевых этапов обработки в живой контейнерной среде приведены в таблице~\ref{tab:paradigm-comparison}.

\begin{table}[htbp]
\centering
\caption{Фактические показатели качества и производительности компонентов в среде исполнения}
\label{tab:paradigm-comparison}
\begin{tabular}{|>{\raggedright\arraybackslash}p{0.34\textwidth}|>{\centering\arraybackslash}p{0.24\textwidth}|>{\raggedright\arraybackslash}p{0.34\textwidth}|}
\hline
\textbf{Компонент / Метрика} & \textbf{Измеренное значение} & \textbf{Режим исполнения в конвейере} \\
\hline
Точность распознавания синтетических партий (7 и 12 полуходов) & \textbf{100\%} (0 нелегальных ходов) & Сквозной прогон с реальной гомографией и полным циклом FSM \\
\hline
Точность репроекции 4 углов доски (\texttt{Hough/RANSAC + CornerSubPix}) & RMSE $\le 2{,}0$~пикс. & Запускается при калибровке и восстановлении геометрии \\
\hline
Инференс \texttt{ChessQ-Lite ViT-S14 (INT8)} & $910{,}98$~мс & Событийно при калибровке и финальной верификации хода \\
\hline
Сравнение патчей \texttt{DINOv2-Small (INT8)} & $145{,}77$~мс & С кэшированием опорного кадра $B_t$ после фиксации хода \\
\hline
Анализ позиции \texttt{Stockfish~16 NNUE (depth=12)} & $115{,}19$~мс & Полностью асинхронно в фоновой очереди \\
\hline
\end{tabular}
\end{table}

Такое разделение частоты вызовов принципиально для работы на CPU без дискретного GPU: тяжёлые трансформерные модели (\texttt{ChessQ-Lite} и \texttt{DINOv2}) вызываются только на ключевых стабилизированных кадрах, благодаря чему потоковый цикл захвата, детекции движения и проверки окклюзии рукой сохраняет отзывчивость в реальном времени.

\textbf{Объяснимость (Explainability)}: в отличие от непрозрачных нейросетевых решений <<чёрного ящика>>, разработанная система формирует количественную декомпозицию для каждого кандидата $m$:
\begin{equation}
S(m) = w_v \cdot S_v(m) + w_t \cdot S_t(m) + w_g \cdot S_g - w_o \cdot P_o(m).
\end{equation}
По каждому полуходу доступны вклады визуальной согласованности $S_v(m)$ (включая оценку зоны основания \textit{Base-Core} и патч-косинусное сходство \texttt{DINOv2}), временной устойчивости $S_t(m)$, геометрической достоверности $S_g$ и штрафа перекрытия рукой $P_o(m)$. При недостаточном запасе уверенности ($\Delta < \tau_{\text{margin}}$) или отсутствии легальных объяснений система воздерживается от искажения протокола, переводя контур в состояние \texttt{UNCERTAIN}.

\subsection{Демонстрация работы гибридной интеллектуальной системы}

Работоспособность системы подтверждена серией сквозных испытаний и интерактивных тестов:

\subsubsection{Сценарий 1: Автоматическая регистрация короткой партии (Детский мат)}
Партия разыгрывалась по тексту \texttt{1. e4 e5 2. Qh5 Nc6 3. Bc4 Nf6 4. Qxf7\# (1-0)}:
\begin{enumerate}
    \item в начале сеанса система локализовала доску, рассчитала матрицу гомографии и перешла из \texttt{IDLE} через \texttt{CALIBRATING} в \texttt{STABLE};
    \item при совершении хода \texttt{1. e4} возмущение перевело FSM в \texttt{MOVING}, а после стабилизации и выдержки дебаунсинга --- в \texttt{VERIFYING};
    \item символьный решатель сопоставил изменённые поля $e2, e4$ со множеством легальных ходов ФИДЕ и подтвердил ход;
    \item на ходе \texttt{4. Qxf7\#} ядро \texttt{python-chess} зафиксировало мат, а \texttt{Stockfish~16} оценил финальную позицию без блокировки интерфейса и без ложной пометки матового хода ошибкой.
\end{enumerate}

\subsubsection{Сценарий 2: Комплекс специальных ходов ФИДЕ}
В сквозном испытании корректно распознана последовательность из 12 полуходов, включающая:
\begin{itemize}
    \item взятие на проходе (\textit{en passant}) с освобождением трёх затронутых полей ($e5, d5, d6$);
    \item короткую рокировку \texttt{O-O} (синхронное перемещение короля $e8 \to g8$ и ладьи $h8 \to f8$);
    \item длинную рокировку \texttt{O-O-O} (перемещение короля $e1 \to c1$ и ладьи $a1 \to d1$);
    \item превращение пешки на 1/8 горизонтали в ферзя.
\end{itemize}

\subsubsection{Сценарий 3: Физическое перекрытие доски рукой без хода}
При задержке руки игрока над центральными полями $d4, d5, e4, e5$ и возврате руки без перестановки фигур детектор рук перевёл FSM в \texttt{MOVING}, затем после очистки сцены и сравнения с опорным кадром $B_t$ визуальные дельты вернулись к нулю, и система безопасно возвратилась в состояние \texttt{STABLE} без ложных записей.

\subsubsection{Сценарий 4: Разрешение неоднозначности оператором (UNCERTAIN)}
При пограничной оценке двух близких легальных кандидатов (например, \texttt{e2e4} и \texttt{e2e3} при малом запасе $\Delta < 0{,}15$) шлюз решений перевёл FSM в \texttt{UNCERTAIN}, отобразил сворачиваемую панель с ранжированными кандидатами, а после подтверждения оператором через веб-панель ход был атомарно зафиксирован.

\subsubsection{Сценарий 5: Обнаружение и регистрация нелегальной перестановки фигур}
При намеренном перемещении фигуры с нарушением правил шахмат (например, прыжок белого коня $g1 \to c5$) гибридный конвейер выявил устойчивое изменение клеток на доске при пустом списке согласованных легальных гипотез. Вместо подбора <<ближайшего>> постороннего хода система:
\begin{enumerate}
    \item перевела автомат в состояние \texttt{UNCERTAIN} без изменения официальной доски $B_t$;
    \item записала событие типа \texttt{illegal\_attempt} в \texttt{events.jsonl} с указанием изменённых полей;
    \item отобразила в веб-панели диагностическое предупреждение с возможностью физического возврата фигур (для автопродолжения) либо ручного ввода допустимого хода UCI.
\end{enumerate}

\subsection{Анализ ограничений и направления развития}

Экспериментальная проверка выявила следующие ограничения прототипа:
\begin{enumerate}
    \item \textbf{Угол наклона камеры}: надёжная детекция основания клеток обеспечивается при угле наклона $50^\circ \le \theta \le 90^\circ$; при меньших углах высокие фигуры перекрывают задние горизонтали.
    \item \textbf{Освещённость и контраст}: при резких точечных бликах и однотонных досках без контраста границ точность геометрической сегментации снижается, поэтому рекомендуется первичная ручная корректировка четырёх углов в веб-панели и отказ от чисто яркостной пороговой сегментации в пользу \textit{Base-Core} и \texttt{DINOv2}.
\end{enumerate}

Направления развития включают адаптацию под одноплатные вычислители (\texttt{Raspberry~Pi~5}, \texttt{Jetson~Orin~Nano}) и расширение дообучения трансформерного классификатора на сложных турнирных наборах деревянных и пластиковых фигур.

\subsection{Выводы по разделу 3}

В результате выполнения третьего этапа курсовой работы:
\begin{enumerate}
    \item Полностью реализован программный комплекс гибридной интеллектуальной системы видеорегистрации шахматных партий, объединяющий алгоритмы компьютерного зрения, квантованные модели \texttt{ChessQ-Lite ViT-S14} и \texttt{DINOv2}, детекцию рук MediaPipe, FSM-дебаунсер, верификатор правил ФИДЕ, транзакционный SHA-256 репозиторий и асинхронный движок \texttt{Stockfish~16}.
    \item Разработан трёхуровневый пользовательский интерфейс: интерактивная веб-панель оператора на \texttt{FastAPI / WebSocket / MJPEG} с калибровкой углов, пошаговым просмотром и диагностикой нелегальных ходов, локальный OpenCV HUD и CLI.
    \item По итогам 6-агентного аудита кодовой базы устранены выявленные дефекты синхронизации, персистентности и локализации; весь комплекс из 285 автоматических тестов выполнен со 100\% успешностью.
    \item На живом контейнере \texttt{chess-recorder-web} подтверждена работоспособность всех контуров и измерены фактические задержки CPU-инференса ключевых моделей.
\end{enumerate}
